\clearpage
\fancyhead[L]{\small CONNAISSANCE DES TEMS.}
\SourceLeaf{614}
\begin{center}\small Göttingische gelehrte Anzeigen. Nos. 40 and 41. Pages 402..408. 11 March 1819.\\[3pt]\rule{.72\linewidth}{.3pt}\end{center}

\noindent\emph{Connaissance des tems ou des mouvemens célestes à l'usage des astronomes et des navigateurs pour l'an {\upshape 1820}, publiée par le bureau des longitudes.} Paris 1818. Published by the widow Courcier. 446 pages in large octavo.

\begingroup\brokenpenalty=10000\relax
In the arrangement of the calendar, we find only one change in this volume. The saints’ names, which oddly enough had been inserted among the astronomical columns for a couple of years, have again been left out. The space thus gained has been used to add a column of differences to the entries for the distance of the First Point of Aries from the Sun and for the Sun’s declination. Something more useful could probably have been chosen in its place. As for the table giving the geographical location of the principal places on Earth, we are assured that it has been checked afresh and that this latest edition is to be regarded as the most accurate and reliable. This assurance, however, is a standing item, and in this volume, as in its predecessors, it takes no long search to see for oneself that the list is edited very carelessly. Thus, for example, we find the latitude of Ofen $32''$ and that of Göttingen $24' 26''$ too high, that of Königsberg $38''$ too low, and that of Mannheim $5''$ too high. For the first and third of these places, points other than the present observatories may have been taken as the basis. This, at least, ought not to have been the case. Among the additions, we touch here only on the original papers, passing over in silence the lengthy extracts from other books and papers, which are very numerous in this volume too. First come the valuable tables for Vesta by \textsc{Daussy}, which are based on the elements and perturbation formulas printed in the volume for 1818. Next comes an instructive paper by \textsc{Laplace} on the apparatus for determining the length of the seconds pendulum. He examines in particular the effect of the knife edge on which the pendulum swings being a portion of a small cylindrical surface instead of a mathematical line in the strict sense. The excellent experiments of \textsc{Borda} on this subject deserved such subtle consideration. --- On determining the perpendicular to the terrestrial meridian and on various other problems relating to it, by \textsc{Puissant}. In our view, the entire approach taken to these subjects so far is still not the right one. Only once they have been considered from new points of view (which has not happened here) can we hope to combine completeness and precision with simplicity and flexibility in the computations. It seems all the more desirable for geometers to turn their attention to this since several governments are currently arranging large-scale surveys which, carried out with the best means available, give reason to hope for new and important insights into the shape of the Earth. According to a report given here by Mr.\ \textsc{Puissant}, the French government has also decided on a new triangulation, which is to be tied to the meridian line of Dunkirk and to the perpendicular from Brest to Strasbourg and extended over the whole kingdom. --- On the elliptical orbit of the comet of 1783 and its resemblance to the orbit of the comet of 1793, by \textsc{Burckhardt}. The poor agreement of the observations of the comet of 1783 with the parabolic elements tried by various astronomers prompted Mr.\ \textsc{Burckhardt} to compute this orbit on the assumption that it was elliptical. Mr.\ \textsc{Burckhardt} obtained an ellipse with a period of $5\tfrac{1}{2}$ years, which represents the observations far better. Even an ellipse with a period of 10 years still gave quite good agreement. Since these elements bear some resemblance to those of the second comet of 1793, it seemed conceivable that the two comets were identical, which is why \SourceLeaf{615}Mr.\ \textsc{Burckhardt} also subjected the observations of the comet of 1793, which had not yet been printed, to a new calculation. Unfortunately, these observations are rather inaccurate. It is scarcely possible to form a definite judgement from the results given here, especially as misprints have probably crept in and their corrections cannot be guessed (the longitude of the node, which \textsc{Saron} had found to be $83^\circ\,55'$ for a parabolic orbit, is given by Mr.\ \textsc{Burckhardt} as $359'\,4'\,48''$ (sic) for one of the elliptical orbits and omitted altogether for a second elliptical orbit). Even so, Mr.\ \textsc{Burckhardt} deserves thanks for drawing astronomers’ attention to this comet and for making the observations of it by \textsc{Messier} known here for the first time. --- Extract from a letter from Mr. \textsc{Schumacher} to Mr. \textsc{Delambre}. The differences between \textsc{Tycho}'s and \textsc{Picard}'s determinations of azimuth on the island of Hueen are satisfactorily explained here on the basis of a paper by \textsc{Augustin}. The Helsingborg tower that \textsc{Picard} had observed was the wrong one, built only after \textsc{Tycho}'s time. Moreover, with imperfect instruments, \textsc{Tycho} had at first laid out his meridian lines with an error of a quarter of a degree, an error which, according to the Historia Coelestis, he himself had already noticed and corrected a few years later. --- On a paper by Mr. \textsc{Littrow} on the thermometric correction of refraction. Mr. \textsc{Delambre} defends at great length here his own table, which Mr. \textsc{Littrow} had criticized as inaccurate. Yet although this defence takes up 10 pages, one still does not find it clearly stated just what Mr. \textsc{Littrow}'s mistake was. Mr. \textsc{Littrow} had hastily assumed that if the density of air, taken as unity at the freezing point, is expressed for a temperature of $\alpha$ degrees above that point by the formula
\[
\frac{1}{1+0.0046875\,\alpha}
\]
 then the \emph{same} formula remains valid when the density at $8^\circ$ is taken as unity and $\alpha$ then denotes the thermometer reading above $8^\circ$. That assumption is plainly incorrect. --- On a new means of regulating the duration of pendulum oscillations, by Mr. \textsc{Prony}. The pendulum rod is extended slightly upward. A very thin, short crossbar is fastened to it by friction so that it can turn in a plane perpendicular to the pendulum rod. The crossbar carries a small platinum sphere at each end. Clearly, the moment of inertia of this pendulum about the axis of oscillation is smallest when the crossbar is parallel to that axis and greatest when the two are at right angles to each other. The oscillations will therefore be fastest in the first case and slowest in the other. With suitably chosen dimensions for the apparatus, this allows very fine adjustment, with the appreciable advantage that the rate can be corrected without stopping the clock. --- On several observations of the planet Uranus found among \textsc{Flamstead}'s stars, by Mr. \textsc{Burckhardt}. Quite by chance, in the course of other searches, Mr. \textsc{Burckhardt} found another \emph{five} observations of Uranus in \textsc{Flamstead}'s Historia Coelestis (1712 March 22; 1715 Feb. 21. 22. 25 and April 18, all Old Style). This discovery is all the more unexpected because all astronomers thought Mr. \textsc{Bode}'s investigation had exhausted the subject, and all the more valuable because the observations from 1715 provide a very good determination of the opposition, whereas the isolated observation from 1690 may still be subject to uncertainty. \textsc{Flamstead} thus observed the planet unknowingly no fewer than six times, and he himself could not possibly have missed the discovery if he had been in the habit of always reviewing his observations immediately on the spot, even if only cursorily, and comparing them with one another. --- On the small equations in Jupiter's motion, by \emph{the very same author}. Mr. \textsc{Burckhardt} gives the numerical values of several equations that \textsc{Laplace} omitted, ranging from the \SourceLeaf{616}second to the sixth order. Some are nevertheless of some significance. Since the corresponding equations omitted in the theory of Saturn's motion must be considerably larger still, it might well be that this circumstance contributed to the striking difference that, as is well known, arose in determining Jupiter's mass from the motion of Saturn and from that of Pallas. --- By \emph{the same author}, remarks on various fixed stars. --- Observations of comets at the Royal Observatory in Paris from 1811---1815 (most already printed elsewhere), together with elements derived from them by Mr. \textsc{Nicollet}. --- Application of the calculus of probabilities to geodetic operations, by \textsc{Laplace}. The calculus of probabilities applies wherever more has been observed than is strictly necessary, provided the observations lack absolute precision. It shows how best to combine the observations and how to determine the relative reliability of the results. Mr. \textsc{Laplace} considers here the case in which two points on the surface of a sphere are joined by a network of triangles, all the angles in the triangles have been observed, and \emph{two} baselines have been measured, one at the beginning and the other at the end of the network. The angle measurements are treated as subject to errors and the baseline as error-free. The quantities sought are the length of the arc joining the two points and the angles it makes with the first and last sides in the network of triangles. It is clear that this is only a very particular problem and that it by no means exhausts the application of probability theory to higher geodesy.
\par\clearpage\endgroup
